% Fragmento público generado desde propietarios íntegros.
% Residencia: 52.13.
% Fuente: ../incoming/radion_monografia_20260827/manuscrito/sections/13_sintesis.tex:1-80
\noindent La synthèse opératoire réunit les lecteurs de phase, d'action, de profondeur et de mémoire sans les identifier entre eux.\par\medskip

\subsubsection{Quadruple réalisation de l'état radional}

Soit \(x_\infty\in X_{\mathrm{enr}}\) la section coinductive. Le radion complet n'est pas
un point isolé, mais le paquet
\begin{equation}
\mathfrak R_\sigma(x_\infty)
=\bigl(
\theta=1,\sigma,
\pih,\alphah,A,C^*,\Delta_4,
\Phi_T,D_E,\epsone,\hret,
q_{\mathrm{mem}},\mathcal I_\infty
\bigr).
\label{eq:full-radion}
\end{equation}
Quatre lecteurs en extraient des objets différents :

\begin{figure}[H]
\centering
\begin{tikzpicture}[
  node distance=17mm and 18mm,
  state/.style={draw=RadNavy,very thick,rounded corners,fill=RadCream,
    minimum width=39mm,minimum height=13mm,align=center},
  readout/.style={draw=RadTeal,rounded corners,fill=RadMist,
    minimum width=39mm,minimum height=12mm,align=center},
  arr/.style={-{Latex},thick,RadCopper}]
\node[state] (x) {état enrichi\\\(x_\infty\)};
\node[readout,above left=of x] (circ) {cercle \(S^1\)\\phase};
\node[readout,above right=of x] (ell) {ellipse positive\\action et forme};
\node[readout,below left=of x] (klein) {interior Klein\\profundidad};
\node[readout,below right=of x] (helix) {hélice nonadique\\mémoire et retour};
\draw[arr] (x)--(circ) node[midway,below left]{\small \(P_{\mathrm{ph}}\)};
\draw[arr] (x)--(ell) node[midway,below right]{\small \(P_{\mathrm{act}}\)};
\draw[arr] (x)--(klein) node[midway,above left]{\small \(P_K\)};
\draw[arr] (x)--(helix) node[midway,above right]{\small \(P_{\mathrm{mem}}\)};
\end{tikzpicture}
\caption{Quatre publications typées du même antécédent. L'unité de
l'état n'élimine pas la différence des codomaines.}
\label{fig:four-readers}
\end{figure}

Le cercle conserve la classe de phase et oublie le relèvement. L'ellipse conserve
l'action, l'anisotropie et l'orientation symplectique. Klein dote l'intérieur d'une
distance projective. L'hélice conserve le nombre de retours et le raffinement.
La mécanique dimensionnelle ne consiste pas à les mélanger, mais à construire les
transducteurs qui les rendent compatibles.

\subsubsection{Interprétation opératoire du cocycle \(\varepsilon_R\)}

\(\varepsilon_R\) n'est pas une constante universelle indépendante de tout contexte. C'est la
coordonnée d'un cocycle de transition associé à la carte du radion. Son
domaine contient deux sections du même état :
\[
S_E=1+D_E,
\qquad
S_T=1+\Phi_T.
\]
Sa règle est
\begin{equation}
\epsone=S_T-S_E=\Phi_T-D_E.
\label{eq:epsilon-final-meaning}
\end{equation}
Son codomaine est la droite adimensionnelle d'unité angulaire ; la carte
\(360/T_+\) la publie en degrés.

\begin{exactbox}
\textbf{En langage simple :} la section directe et la section torsionnelle arrivent
à la même cellule d'unité par deux chemins internes différents. Leurs parties
entières coïncident ; leurs restes, non. \(\epsR\) est la différence orientée de ces
restes, calculée avec la retenue d'APP et portée à la limite par les cylindres
TPK. Elle n'est pas introduite pour que le radian se ferme : elle apparaît avant de consulter le
radian et, lorsqu'elle est publiée, elle le ferme.
\end{exactbox}

Cette définition explique pourquoi le nombre est petit. Non parce qu'on a
cherché une petite correction, mais parce que \(S_E\) et \(S_T\) partagent un
préfixe profond. La petitesse mesure la proximité structurelle de deux
publications ; le signe mesure leur orientation ; les chiffres enregistrent la retenue.

